\documentclass[lettersize,journal]{IEEEtran}
\usepackage{amsmath,amsthm,amssymb}
\usepackage{mathrsfs}
\usepackage{amsfonts}
%\usepackage{algorithmic}
\usepackage[ruled,vlined,linesnumbered]{algorithm2e}
% 自定义算法宽度

%\usepackage{algorithmicx}
%\usepackage{algorithm}
\usepackage{array}
\usepackage{makecell}
\usepackage{cite}
\usepackage{textcomp}
\usepackage{stfloats}
\usepackage{bbding}

%————————————————————————————————————————
 \usepackage{etoolbox}
% \usepackage{ulem}
 \usepackage{lipsum} 
 \usepackage{tikz}
 \usepackage{enumitem}
 \usepackage[most]{tcolorbox} % 引入 tcolorbox 宏包

\tcbset{
  mybox/.style={
    breakable,           % 允许跨页
    enhanced jigsaw,     % 增强拼图模式，允许更复杂的内容布局
    colframe=black,      % 边框颜色
    colback=white,       % 背景颜色
    boxrule=0.2mm,       % 边框厚度
    width=\columnwidth,    % 框的宽度
    sharp corners,       % 直角边框
    before skip=10pt,    % 框上方的间距
    after skip=10pt,     % 框下方的间距
    colbacktitle=white,  % 标题背景颜色
    title=\quad,               %页和跨栏的框,
    attach title to upper, % 标题附加到框的上部
    fonttitle=\bfseries, % 标题字体样式
    center title,        % 标题居中
    boxsep=5pt,          % 框内内容与边框的距离
    left=4pt,            % 左边距
    right=4pt,           % 右边距
    top=-9pt,            % 上边距
    bottom=-1pt          % 上边距
  }
}

\usepackage{graphicx}
\makeatletter
\newcommand*\bigcdot{\mathpalette\bigcdot@{.5}}
\newcommand*\bigcdot@[2]{\mathbin{\vcenter{\hbox{\scalebox{#2}{$\m@th#1\bullet$}}}}}
\makeatother
%----------------------------------------

\usepackage{url}
\usepackage{color}
\usepackage{xcolor}
\usepackage[hidelinks]{hyperref}
\usepackage[referable]{threeparttablex}


\usepackage{multirow} 
\usepackage{pifont}

\usepackage{subfigure}
\newtheorem{definition}{Definition}
\newtheorem{theorem}{Theorem}
\newtheorem{lemma}{Lemma}
\newtheorem{remark}{Remark}
\usepackage{balance}
\usepackage{orcidlink}
\usepackage{pifont}
\newcommand{\circled}[1]{{\normalsize{\textcircled{\scriptsize{#1}}}\normalsize\;}}


\SetAlCapSkip{0.5em} % 减少算法标题与内容的间距

\setlength{\algomargin}{1em} % 调整算法整体的左右间距

\newcommand{\code}[1]{\texttt{#1}} %%% 


\newcommand{\crossmeta}[0]{\ensuremath{\code{MetaBridge}\;}}


%% ---- Comments Format, added by Huawei
\newcommand{\hhw}[1]{\textcolor{black}{{#1}}} %COMMENTS BY Huawei magenta

\newcommand{\hw}[1]{\textcolor{black}{{#1}}} %COMMENTS BY Huawei magenta


\newcommand{\YCL}[1]{\textcolor{black}{{#1}}} %COMMENTS BY Changlin gray
\newcommand{\Yng}[1]{{\color{black}#1}}%COMMENTS BY Qinglin blue
\newcommand{\yzk}[1]{{\color{black}#1}} %COMMENTS BY Zhaokang red
\newcommand{\ltt}[1]{\textcolor{black}{{#1}}} %COMMENTS BY LI Taotao brown
\newcommand{\lt}[1]{\textcolor{black}{{#1}}} %COMMENTS BY LI Taotao brown
\newcommand{\lxf}[1]{\textcolor{black}{{#1}}} %COMMENTS BY Xiaofei brown
%% ---- Comments Format:~
\newcommand{\revised}[1]{\textcolor{magenta}{{#1}}}
\newcommand{\todo}[1]{\textcolor{red}{{#1}}} 


\hyphenation{op-tical net-works semi-conduc-tor IEEE-Xplore}


\begin{document}
\title{\code{FastOracle}: Improving Oracle Throughput and Latency via a Reduced-Size Secure Committee}


% \author{Taotao~Li,~Huawei~Huang,~\IEEEmembership{Senior Member,~IEEE,} Xiaofei~Luo, Qinglin~Yang,~\IEEEmembership{Member,~IEEE,}~Zibin~Zheng,~\IEEEmembership{Fellow,~IEEE}
% %Taotao~Li~\orcidlink{0000-0002-5348-442X},
% %Huawei~Huang~\orcidlink{0000-0002-7035-6446},
% %Xiaofei~Luo~\orcidlink{0000-0001-8258-852X},
% %Qinling~Yang~\orcidlink{0000-0002-7263-8914},
% %Anke~Xie,
% %Zibin~Zheng~\orcidlink{0000-0002-7878-4330}\\
    

% \thanks{Taotao Li, Huawei Huang Xiaofei Luo and Zibin Zheng are with the School of Software Engineering, Sun Yat-Sen University, Zhuhai, 528478, China. Taotao Li is also with the Yunnan Innovation Institute of Beihang University, Yunnan Key Laboratory of Blockchain Application Technology, Kunming, 650233, China. (E-mail:\{litt93; huanghw28; luoxf37; zhzibin\}@mail.sysu.edu.cn)}
% %
% \thanks{Qinglin Yang is with Cyberspace Institute of Advanced Technology, Guangzhou University, Guangzhou, 510006, China. (E-mail: yangqinglin@gzhu.edu.cn)} 


% %\thanks{Debiao He is with the School of Cyber Science and Engineering, Wuhan University, Wuhan 430072, China. (E-mail: hedebiao@163.com)}
% %\thanks{Yan Zhang is with the Department of Informatics, University of Oslo, Norway. (Email: yanzhang@ieee.org)}
% \thanks{Corresponding author: Huawei Huang.}

% }




\markboth{Journal of \LaTeX\ Class Files,~Vol.~14, No.~8, August~2021}%
{Shell \MakeLowercase{\textit{et al.}}: A Sample Article Using IEEEtran.cls for IEEE Journals}


\maketitle


\begin{abstract}

Oracle

\end{abstract}


\begin{IEEEkeywords}
Oracle, Blockchain, Committee, Throughput.
\end{IEEEkeywords}


\section{Introduction}
\label{sec:intro}

\IEEEPARstart{O}{racle}  

\section{SYSTEM DESIGN}

\subsection{Committee size}
\label{sec:calculate_committee_size}

Let \(P_a\) denote the maximum ratio of adversarial population in the pre-selected pool. Since committee members are sampled uniformly from this pool, the probability that any given committee position is occupied by a malicious node remains \(P_a\).  A committee of size $N_c$ is considered compromised if the adversarial node controls strictly more than $\lfloor N_c / 2 \rfloor$ of its members, as the \code{FastOracle} consensus protocol requires a strict majority (greater than 1/2) of nodes to be honest. Letting $X$ denote the number of malicious nodes in a committee, the probability $P_f$ that a certain committee is compromised is:

\begin{align}
P_f 
&= P\!\left(X > \left\lfloor\frac{1}{2}N_c \right \rfloor\right)\notag \\
&= \sum_{i= \left\lfloor\frac{1}{2}N_c \right\rfloor+ 1}^{N_c} P(X=i) \notag \\
&= \sum_{i= \left\lfloor\frac{1}{2}N_c \right\rfloor+ 1}^{N_c} \binom{N_c}{i} P_a^{i} (1-P_a)^{N_c-i}
\end{align}

In our proposed \code{FastOracle} protocol, security requires that all parallel committees behave honestly. Let $N$ be the number of parallel committees and $N_{\text{pre}}$ the size of the pre-selected pool. The probability $P_{\text{all}}$ that at least one committee is compromised yields the equation:

\begin{equation}
P_{\text{all}}=1-(1-P_f)^N=1-(1-P_f)^{\lfloor N_{\text{pre}} / N_c \rfloor}
\end{equation}

Note that the process of selecting the pre-selected committee from the underlying blockchain network fails with probability \(2^{-\beta}\). Let \(P\) denote the end-to-end probability that our \code{FastOracle} is compromised. Then we have the following equation:

\begin{align}
P &= 1 - \left(1 - \frac{1}{2^{\beta}}\right)(1 - P_{\text{all}}) \notag\\
  &= 1 - \left(1 - \frac{1}{2^{\beta}}\right)(1 - P_f)^{\lfloor N_{\text{pre}} / N_c \rfloor} \notag\\
&= 1 - \left(1 - \frac{1}{2^{\beta}}\right) \notag\\
  &\quad\times\Biggl(1 - \sum_{i= \left\lfloor\frac{1}{2}N_c \right\rfloor+ 1}^{N_c} \binom{N_c}{i} P_a^{i} (1-P_a)^{N_c-i} \Biggr)^{\lfloor N_{\text{pre}} / N_c \rfloor}
\end{align}

To achieve a predefined end-to-end security level of $\alpha$ bits (i.e., $P < 2^{-\alpha}$), for fixed security parameter $\beta$, adversarial ratio \(P_a\) and pre-selected pool size $N_{\text{pre}}$, we select the smallest committee size $N_c$ such that the inequality \(P < 2^{-\alpha}\) holds. This yields the maximal number of parallel committees, thereby maximizing the \code{FastOracle}  throughput, while ensuring that the probability of \code{FastOracle} being compromised remains below \(2^{-\alpha}\).

\section{SECURITY ANALYSIS}

In this section, we analyze \code{FastOracle} from Section
? satisfies the security properties defined in Section ?.

\subsection{Threat Model}

\section{IMPLEMENTATION AND EVALUATION}

Next, we implement the proposed \code{FastOracle}, and evaluate its performance in terms of request processing time, request processing throughput, queue size of request pool, and the number of generated certificates. 

\subsection{Experimental setup}
\textbf{Basic Settings.} We conduct our \code{FastOracle} prototype in
Golang and evaluated its performance on Proof-of-Stake (PoS) blockchain and Proof-of-Work (PoW) blockchain, respectively. The prototype was deployed on a Windows system with an Intel CPU Core i7-12700H (2.30GHz) and 32GB of DDR5 RAM. 

All experiments were conducted with a oracle network size of \(N=2000\) nodes. Following the committee-size analysis in Section \ref{sec:calculate_committee_size}, \code{FastOracle} organizes nodes into four parallel committees, each comprising 400 nodes. The maximum number of Byzantine nodes was fixed to \(\lfloor \frac{N-1}{3} \rfloor\).

Each experiment consisted of 20,000 oracle requests injected at a sustained rate of 100 requests per second. Requests were processed in batches of 15 to maximize throughput while preserving liveness. To reflect real-world wide-area network conditions, inter-node message latency was drawn from a uniform distribution with a maximum of 2 seconds and an expected value of 1.5 seconds. External data were fetched via HTTPS from the Shenzhen Stock Exchange under a constrained downlink bandwidth of 100 Mb/s.

\textbf{Chain Settings.} We deployed an Ethereum-compatible blockchain using Ganache. To emulate representative PoS and PoW environments, we configured average block intervals of 12 seconds (Ethereum-like PoS) and 600 seconds (Bitcoin-like PoW), respectively.

\textbf{Baseline Solutions.} For comprehensive benchmarking, we also implemented four state-of-the-art oracle schemes in Golang:

\begin{itemize}
    \item \textbf{SenFeed}: A dynamic decentralized oracle for IoT streaming data. It employs incremental truth discovery for real-time accurate inference and a novel BFT consensus protocol supporting dynamic node membership without system resets.
    
    \item \textbf{DeepThought}: A human-in-the-loop oracle that certifies off-chain information via distributed voting and reputation mechanisms, achieving stronger resistance to voter corruption than prior reputation-based systems.
    
    \item \textbf{DECENTRUTH}: A truth-discovery oracle addressing deceitful but authentic sources. It combines off-chain truth discovery with asynchronous BFT consensus to deliver highly accurate estimated truths.
    
    \item \textbf{DAON}: A fully autonomous oracle network with a dedicated consensus protocol, non-interactive reputation and payment schemes, and continuous security monitoring, ensuring tamper-proof data feeds for smart contracts in complex Byzantine environments.
\end{itemize}

These baselines enable a fair and reproducible comparison with \code{FastOracle} under identical hardware, network, and workload conditions.

\textbf{Metrics.} We evaluate our \code{FastOracle} using the following four metrics: \textbf{i) Request processing time.} The time required to process a single request, measured from the moment the request is submitted to the oracle network until the requested data is successfully written on-chain. \textbf{ii) Request processing throughput.} The number of requests processed per second by oracle system. \textbf{iii) Queue size of request pool.} The number of pending requests in the oracle request pool. \textbf{iv) Number of generated certificates.} The number of authenticity attestations generated by the oracle nodes per second. 

 

% Figure for POS Network
\begin{figure}[!t]
    \centering
    % Make sure the path matches your folder structure exactly
    \includegraphics[width=\linewidth]{figure/figs_pos_basic3/stacked_time_bars.pdf}
    \caption{Latency breakdown of different oracle protocols in the PoS network environment. The total processing time is dissected into three phases: search, consensus, and on-chain.}
    \label{fig:pos_latency}
\end{figure}

% Figure for POW Network
\begin{figure}[!t]
    \centering
    % Make sure the path matches your folder structure exactly
    \includegraphics[width=\linewidth]{figure/figs_pow_basic3/stacked_time_bars.pdf}
    \caption{Latency breakdown of different oracle protocols in the PoW network environment. A broken y-axis is utilized to accommodate the significant disparity in total processing time, specifically highlighting the extensive on-chain latency of Deepthought.}
    \label{fig:pow_latency}
\end{figure}

  

% 1. Consensus Time Stability for POS (Standard View)
\begin{figure}[!t]
    \centering
    % Check if the path 'figure/figs_pos_basic3/' matches your project structure
    \includegraphics[width=\linewidth]{figure/figs_pos_basic3/consensusTime_stability.pdf}
    \caption{Stability analysis of consensus time in the PoS network. The solid line represents the rolling mean (window=50), while the shaded region indicates the standard deviation. The results demonstrate the low-latency consistency of the proposed protocol compared to baselines.}
    \label{fig:pos_consensus_stability}
\end{figure}

% 2. Consensus Time Stability for POW (Broken Axis View)
\begin{figure}[!t]
    \centering
    % Check if the path 'figure/figs_pow_basic3/' matches your project structure
    \includegraphics[width=\linewidth]{figure/figs_pow_basic3/consensusTime_stability.pdf}
    \caption{Stability analysis of consensus time in the PoW network. A broken y-axis is employed to simultaneously visualize the high-latency Deepthought protocol ($\approx$600s) and the lower-latency group. The narrow variance band of our protocol (Ours) indicates robust performance stability even under high-difficulty PoW conditions.}
    \label{fig:pow_consensus_stability}
\end{figure}



\end{document}